\section{Schallgeschwindigkeit in Festkörpern}

\begin{aufgabe}{Rohdaten}
  Stellen Sie die gemessenen Werte von $L$, $D$ und $M$ in
  tabellarischer Form dar. Visualisieren Sie für jeden vermessenen
  Metallstab einen typischen Schwingungsverlauf sowie dessen
  Fourierspektrum in geeigneter Weise.
\end{aufgabe}

% Bitte belassen Sie die Aufgabentexte in Ihrem Protokoll und beginnen
% Sie hier mit der Lösung der ersten Aufgabe:

% Achten Sie darauf, alle verwendeten Formeln anzugeben und die
% verwendeten Formelzeichen zu definieren.

\begin{aufgabe}{Auswertung}
  Bestimmen Sie aus den aufgezeichneten Schwingungsvorgängen die
  Schwingungsfrequenzen und tabellieren Sie die erhaltenen
  Werte. Ermitteln Sie die Streuung der Frequenzwerte und bestimmen
  Sie daraus den statistischen Fehler auf den Mittelwert. Erläutern
  und illustrieren Sie Ihr Vorgehen zur Bestimmung der systematischen
  Unsicherheiten auf die Frequenz. Diskutieren Sie die Unsicherheiten
  auf $L$, $D$ und $M$. Berechnen Sie den E-Modul $E$ und die Dichte
  $\rho$ der vorliegenden Stangen und die zugehörigen statistischen
  und systematischen Messunsicherheiten. Fassen Sie Ihre Ergebnisse
  tabellarisch zusammen. Diskutieren Sie, welche Fehlerbeiträge den
  Gesamtfehler auf $E$ dominieren. Vergleichen Sie Ihre Ergebnisse für
  $E$ und $\rho$ mit einschlägigen Literaturwerten.
\end{aufgabe}



